\documentclass[11pt,a4paper]{../_shared/altacv}

\geometry{left=1.5cm,right=1.5cm,top=1.cm,bottom=1.2cm,columnsep=1.5cm}

\usepackage[utf8]{inputenc}
\usepackage[T1]{fontenc}
\usepackage[french]{babel}
\usepackage{paracol}
\usepackage{xcolor}
\usepackage{hyperref}
\usepackage{fontawesome5}
\usepackage{setspace}

\definecolor{SlateGrey}{HTML}{2E2E2E}
\definecolor{LightGrey}{HTML}{666666}
\definecolor{DarkPastelRed}{HTML}{8AC0D9}
\definecolor{PastelRed}{HTML}{00B0DA}
\definecolor{GoldenEarth}{HTML}{B4AD0C}
\colorlet{name}{black}
\colorlet{tagline}{PastelRed}
\colorlet{heading}{DarkPastelRed}
\colorlet{headingrule}{GoldenEarth}
\colorlet{subheading}{PastelRed}
\colorlet{accent}{PastelRed}
\colorlet{emphasis}{SlateGrey}
\colorlet{body}{LightGrey}

\renewcommand{\familydefault}{\sfdefault}

\name{\Large BOILEAU Guillaume}
\personalinfo{
  \email{guillaume.boileaupro@gmail.com}
  \phone{+33 6 59 94 85 84}
  \location{Alpes Maritimes, France}
  \linkedin{guillaume-boileau-phd-891b81265}
  \researchgate{Guillaume-Boileau-2}
  \orcid{0000-0002-3576-6968}
}

\setlength{\parskip}{0.75\baselineskip}

\begin{document}
\makecvheader
\vspace{-0.6cm}
\cvsection{Lettre de motivation}
\vspace{-0.4cm}

{\color{accent}\textbf{Objet :}} \textit{Candidature au poste de Chargé d'études et de projets H/F (31052)}

{\color{body}

Madame, Monsieur,

Je vous adresse ma candidature au poste de Chargé d'études et de projets au sein de la Direction de la construction, de l'immobilier et du patrimoine. Ce poste retient particulièrement mon attention car il associe plusieurs dimensions qui correspondent directement à mon parcours et à ma façon de travailler : conduite de projets techniques, suivi de systèmes informatiques et d'équipements, coordination d'intervenants, structuration de déploiements, contrôle de la bonne exécution et amélioration continue de dispositifs opérationnels.

Mon parcours m'a conduit à évoluer dans des environnements techniques complexes, où la qualité de préparation, la rigueur d'organisation, la fiabilité des systèmes et la capacité à sécuriser l'exécution étaient des exigences permanentes. J'y ai développé une compétence solide en gestion de projet technique, avec une implication concrète sur l'ensemble du cycle de vie des sujets confiés : analyse des besoins, cadrage, planification, priorisation des actions, coordination des contributions, suivi de l'avancement, gestion des points de blocage, mise en oeuvre opérationnelle et contrôle des résultats. Cette culture du pilotage constitue aujourd'hui l'un des axes forts de mon profil.

Dans mes expériences précédentes, j'ai eu à structurer des activités techniques impliquant des contraintes multiples, des interlocuteurs variés et des objectifs élevés de qualité et de robustesse. J'ai assuré l'organisation de travaux, la répartition des priorités, le suivi de livrables, l'encadrement de contributeurs, la coordination entre besoins fonctionnels et solutions techniques, ainsi que le maintien d'une vision d'ensemble permettant de tenir les délais et de sécuriser l'avancement. Cette capacité à faire converger des actions techniques vers un résultat clair, exploitable et durable me paraît pleinement en phase avec les attendus du poste.

J'ai également développé une solide base informatique, en cohérence avec vos prérequis. J'ai travaillé sur des systèmes logiciels nécessitant fiabilité, supervision, continuité de fonctionnement et capacité de diagnostic. J'ai participé au développement, à l'intégration et au suivi de services techniques dans des contextes exigeants, avec une attention constante portée à la robustesse des traitements, à la qualité d'exécution, à la traçabilité des opérations et à la résolution rapide des incidents. Cette expérience me permet d'aborder avec sérieux des missions impliquant le suivi de systèmes informatiques de sûreté, l'analyse de dysfonctionnements, l'amélioration du fonctionnement global et la contribution à des dispositifs techniques sensibles.

Mon expérience récente a renforcé cette dimension très concrète de gestion de systèmes techniques. J'y ai travaillé sur des services et des flux liés à des équipements physiques, dans un contexte où il fallait assurer à la fois l'intégration, le bon fonctionnement opérationnel, le suivi d'interfaces techniques, la résolution d'incidents et la continuité de service. Cette expérience m'a appris à intervenir avec méthode sur des sujets où la technique ne peut pas rester théorique : il faut comprendre le système, identifier les causes d'un dysfonctionnement, coordonner les actions correctives, suivre les mises en oeuvre et garantir que la solution fonctionne réellement dans la durée. Cette logique me semble directement transposable à la gestion de projets liés au contrôle d'accès, à la vidéoprotection ou aux systèmes d'alarme.

Au-delà de la seule dimension informatique, je peux apporter une vraie valeur sur la partie gestion de projet. J'ai l'habitude de piloter des sujets techniques de manière structurée, avec une approche fondée sur l'anticipation, la méthode et le suivi. Cela recouvre concrètement la définition des objectifs, l'identification des contraintes, l'organisation du travail, la formalisation des besoins, le suivi de l'exécution, la coordination des intervenants, le reporting d'avancement, l'identification des risques, la mise en place d'actions correctives et la vérification de la conformité du résultat livré. J'accorde également une grande importance à la lisibilité des projets, à la qualité de la documentation, à la traçabilité des décisions et à la capacité à produire des outils de suivi utiles à la décision.

J'ai aussi une expérience réelle du travail transverse, qui me paraît essentielle pour ce poste. Dans mes fonctions précédentes, j'ai dû travailler à l'interface entre plusieurs contributeurs, plusieurs niveaux de responsabilité et plusieurs contraintes parfois concurrentes. J'ai appris à faire dialoguer besoin, technique, exécution et calendrier, en maintenant un cadre clair pour faire avancer les sujets. Cette capacité relationnelle et organisationnelle est, selon moi, indispensable pour suivre des prestataires, contribuer à l'élaboration de dossiers techniques, organiser des actions de maintenance, contrôler la bonne réalisation des opérations et garantir la cohérence d'ensemble des projets déployés.

Le poste proposé m'intéresse aussi par sa dimension de terrain et de service. J'apprécie les fonctions où il faut articuler réflexion technique, pilotage concret, observation des installations, diagnostics, suivi de dispositifs et amélioration continue. Travailler sur des équipements ayant un impact direct sur la sécurité des personnes et des biens donne un sens immédiat à l'action. Je suis particulièrement motivé par l'idée de mettre mes compétences en organisation, en informatique et en gestion de projets techniques au service d'une collectivité, dans un cadre où la rigueur, la fiabilité et l'efficacité ont une portée très concrète.

Autonome, rigoureux et impliqué, je suis à l'aise dans les contextes qui demandent à la fois capacité d'analyse, sens des responsabilités, prise de décision, travail en équipe et suivi opérationnel dans la durée. Je suis convaincu que mon expérience de pilotage technique, ma capacité à structurer des projets complexes, ma maîtrise des environnements informatiques et mon goût pour les sujets concrets peuvent constituer un apport réel pour votre direction.

Je serais heureux de pouvoir échanger avec vous afin de vous présenter plus en détail mon parcours et ma motivation.

Je vous prie d'agréer, Madame, Monsieur, l'expression de mes salutations distinguées.

\textbf{Guillaume Boileau}

}

\end{document}